\documentclass{../templateReading/cdreading}

\newcommand{\authorinfo}{THE FORUM CENTER - NGUYỄN HOÀNG HUY}
\newcommand{\documenttitle}{TÀI LIỆU CHUYÊN ĐỀ READING}
\newcommand{\collectiontitle}{Level 1 — W5 (Jan 26-Feb 1)}

\begin{document}

\thispagestyle{firstpage}
\makeworkshopheader{\authorinfo}{\documenttitle}{\collectiontitle}
\pagestyle{otherpages}

\cdsection{The Effects of Deforestation}

\textbf{A}\par

Every year it is estimated that roughly 5.2 million hectares (52,000 km2 of the forest is lost worldwide. That is a net figure, meaning it represents the area of the forest not replaced. T put this size in context, that is an area of land the size of Croatia lost every single year. There is a wide range of negative effects from deforestation that range from the smallest biological processes right up to the health of our planet as a whole. On a human level, millions of lives are affected every year by flooding and landslides that often result from deforestation.

\textbf{B}\par

There are 5 million people living in areas deemed at risk of flooding in England and Wales. Global warming, in part, worsened by deforestation, is responsible for higher rainfalls in Britain in recent decades. Although it can be argued that demand for cheap housing has meant more houses are being built in at-risk areas, the extent of the flooding is increasing. The presence of forests and trees along streams and rivers acts like a net. The trees catch and store water, but also hold the soil together, preventing erosion. By removing the trees, the land is more easily eroded increasing the risk of landslides and also, after precipitation, less water is intercepted when trees are absent and so more enters rivers, increasing the risk of flooding.

\textbf{C}\par

It is well documented that forests are essential to the atmospheric balance of our planet, and therefore our own wellbeing too. Scientists agree unequivocally that global warming is a real and serious threat to our planet. Deforestation releases 15\% of all greenhouse gas emissions. One-third of the carbon dioxide emissions created by human activity comes from deforestation around the globe.

\textbf{D}\par

In his book Collapse, about the disappearance of various ancient civilisations, writer Jared Diamond theorises about the decline of the natives of Easter Island. European missionaries first arrived on the island in 1722. Research suggested that the island, whose population was in the region of two to three thousand at the time, had once been much higher at fifteen thousand people. This small native population survived on the island despite there being no trees at all. Archaeological digs uncovered evidence of trees once flourishing on the island. The uncontrolled deforestation not only led to the eradication of all such natural resources from the island but also greatly impacted the number of people the island could sustain. This underlines the importance of forest management, not only for useful building materials but also for food as well.

\textbf{E}\par

Forestry management is important to make sure that stocks are not depleted and that whatever is cut down is replaced. Without sustainable development of forests, the levels of deforestation are only going to worsen as the global population continues to rise, creating a higher demand for the products of forests. Just as important though is consumer awareness. Simple changes in consumer activity can make a huge difference. These changes in behaviour include, but are not limited to, recycling all recyclable material; buying recycled products and looking for the FSC sustainably sourced forest products logo on any wood or paper products.

\textbf{F}\par

Japan is often used as a model of exemplary forest management. During the Edo period between 1603 and 1868 drastic action was taken to reverse the country’s serious exploitative deforestation problem. Whilst the solution was quite complex, one key aspect of its success was the encouragement of cooperation between villagers. This process of collaboration and re-education of the population saved Japan’s forests. According to the World Bank, 68.5\% of Japanese land area is covered by forest, making it one of the best performing economically developed nations in this regard.

\textbf{G}\par

There is, of course, a negative impact of Japan’s forest management. There is still a high demand for wood products in the country, and the majority of these resources are simply imported from other, poorer nations. Indonesia is a prime example of a country that has lost large swaths of its forest cover due to foreign demand from countries like Japan. This is in addition to other issues such as poor domestic forest management, weaker laws and local corruption. Located around the Equator, Indonesia has an ideal climate for the rainforest. Sadly much of this natural resource is lost every year. Forest cover is now down to less than 51\% from 65.4\% in 1990. This alone is proof that more needs to be done globally to manage forests.

\textbf{H}\par

China is leading the way in recent years for replenishing their forests. The Chinese government began the Three-North Shelter Forest Program in 1978, with aims to complete the planting of a green wall, measuring 2,800 miles in length by its completion in 2050. Of course, this program is in many ways forced by nature itself; the expansion of the Gobi Desert threatened to destroy thousands of square miles of grassland annually through desertification. This is a process often exacerbated by deforestation in the first place, and so represents an attempt to buck the trend. Forested land in China rose from 17\% to 22\% from 1990 to 2015 making China one of the few developing nations to reverse the negative trend.

\cdsection{Glossary}

exemplary: serving as a perfect example

exacerbate: make worse

\cdsection{Questions 14-20}

The reading passage below has eight paragraphs, A-H.

\cdnoteinline{Choose the correct heading for each paragraph from the list of headings below.}

\cdnoteinline{Write the correct number, i-xi, in boxes 14-20 on your answer sheet.}

\cdblue{\textbf{i}} Atmospheric impacts

\cdblue{\textbf{ii}} Ideal forestry management example

\cdblue{\textbf{iii}} No trees, less people

\cdblue{\textbf{iv}} Good uses for wood

\cdblue{\textbf{v}} Looking after the forests

\cdblue{\textbf{vi}} Numbers of lost trees

\cdblue{\textbf{vii}} Wasted water

\cdblue{\textbf{viii}} Replanting forests

\cdblue{\textbf{ix}} Happy trees

\cdblue{\textbf{x}} Flood risks

\cdblue{\textbf{xi}} Poorer nations at higher risk

Example             Answer

Paragraph A      vi

\textbf{14} Paragraph B

\textbf{15} Paragraph C

\textbf{16} Paragraph D

\textbf{17} Paragraph E

\textbf{18} Paragraph F

\textbf{19} Paragraph G

\textbf{20} Paragraph H

\cdsection{Questions 21-26}

\cdnoteinline{Complete the summary below.}

\cdnoteinline{Choose NO MORE THAN TWO WORDS AND/OR A NUMBER from the passage for each answer.}

\cdnoteinline{Write your answers in boxes 21-26 on your answer sheet.}

The effects of deforestation are widespread and various. Some examples include flooding at a local scale to the wider effects of global warming on a worldwide scale. In Britain, for example, 21 …………………. people live in areas at risk of flooding. This risk is increased by deforestation. Trees catch and 22 ………………… water lowering the chance of flooding. By removing trees land erosion is also higher, increasing the chance of 23 ………………… Deforestation also affects global warming by contributing 15\% of the 24 …………………. of greenhouse gasses. To make sure that the cutting down of trees is done in a sustainable way, good forestry 25 …………………. is important. In most countries, more trees are cut down every year than planted. One country that is reversing this trend is China, making it one of the few nations to 26 ………………….. the more common negative trend.

\cdblue{\textbf{Hướng dẫn giải}}

\cdsection{Phần I: Câu hỏi 14-20 (Matching Headings)}

Mục tiêu: Nắm bắt ý chính của đoạn văn và tìm tiêu đề tương ứng trong danh sách cho sẵn.

\cdstep{1. Chiến lược tiếp cận}
\begin{itemize}
\item Bước 1: Giải mã Danh sách Tiêu đề (List of Headings)
\begin{itemize}
\item Đọc kỹ danh sách các tiêu đề từ (i) đến (xi).
\item Gạch chân từ khóa và tóm lược nội dung chính:
\begin{itemize}
\item (i) Atmospheric impacts (Tác động đến khí quyển).
\item (ii) Ideal forestry management example (Ví dụ quản lý rừng lý tưởng).
\item (iii) No trees, less people (Không có cây, ít người hơn).
\item (iv) Good uses for wood (Công dụng tốt của gỗ).
\item (v) Looking after the forests (Chăm sóc rừng).
\item (vi) Numbers of lost trees (Số lượng cây bị mất).
\item (vii) Wasted water (Nước bị lãng phí).
\item (viii) Replanting forests (Trồng lại rừng).
\item (ix) Happy trees (Cây hạnh phúc).
\item (x) Flood risks (Nguy cơ lũ lụt).
\item (xi) Poorer nations at higher risk (Các quốc gia nghèo hơn chịu rủi ro cao hơn).
\end{itemize}
\end{itemize}
\item Bước 2: Đọc lướt (Skimming) và Đối chiếu
\begin{itemize}
\item Đọc lướt từng đoạn văn để tìm câu chủ đề hoặc ý bao trùm.
\item Lưu ý: Tìm sự tương đồng về ý nghĩa giữa tiêu đề và nội dung đoạn văn.
\end{itemize}
\end{itemize}

\cdstep{2. Phân tích chi tiết từng đoạn văn}
\begin{itemize}
\item Câu 14: Paragraph B
\begin{itemize}
\item Nội dung chính: Đoạn văn mở đầu bằng việc nhắc đến "5 million people living in areas deemed at risk of flooding". Tiếp theo giải thích việc chặt cây làm tăng nguy cơ này: "increasing the risk of flooding".
\item Kết luận: Nội dung tập trung hoàn toàn vào rủi ro ngập lụt, khớp với tiêu đề (x) Flood risks.
\item Đáp án: x
\end{itemize}
\item Câu 15: Paragraph C
\begin{itemize}
\item Nội dung chính: Đoạn văn đề cập đến "atmospheric balance" (cân bằng khí quyển), "global warming" (nóng lên toàn cầu), và "greenhouse gas emissions" (khí thải nhà kính).
\item Kết luận: Các từ khóa này đều liên quan đến tác động lên không khí/khí quyển, khớp với tiêu đề (i) Atmospheric impacts.
\item Đáp án: i
\end{itemize}
\item Câu 16: Paragraph D
\begin{itemize}
\item Nội dung chính: Đoạn văn kể về sự suy tàn của người bản địa đảo Easter. Tác giả viết: "The uncontrolled deforestation... greatly impacted the number of people the island could sustain." (Nạn phá rừng không kiểm soát... ảnh hưởng lớn đến số lượng người mà hòn đảo có thể nuôi sống).
\item Kết luận: Ý chính là việc mất rừng dẫn đến sự suy giảm dân số, khớp với tiêu đề (iii) No trees, less people.
\item Đáp án: iii
\end{itemize}
\item Câu 17: Paragraph E
\begin{itemize}
\item Nội dung chính: Đoạn văn bàn về "Forestry management" (quản lý lâm nghiệp) để đảm bảo bền vững. Nó cũng nhắc đến "consumer awareness" (nhận thức người tiêu dùng) và các hành động bảo vệ như "recycling" (tái chế), "sustainably sourced" (nguồn gốc bền vững).
\item Kết luận: Tất cả các hành động này đều nhằm mục đích bảo vệ và chăm sóc rừng, khớp với tiêu đề (v) Looking after the forests.
\item Đáp án: v
\end{itemize}
\item Câu 18: Paragraph F
\begin{itemize}
\item Nội dung chính: Câu đầu tiên khẳng định: "Japan is often used as a model of exemplary forest management" (Nhật Bản thường được dùng làm mô hình quản lý rừng kiểu mẫu).
\item Kết luận: Từ "exemplary" (kiểu mẫu) đồng nghĩa với "ideal" (lý tưởng), khớp với tiêu đề (ii) Ideal forestry management example.
\item Đáp án: ii
\end{itemize}
\item Câu 19: Paragraph G
\begin{itemize}
\item Nội dung chính: Đoạn văn nói về tác động tiêu cực của việc Nhật Bản nhập khẩu gỗ từ "poorer nations" (các quốc gia nghèo hơn) như Indonesia. Indonesia bị mất rừng do nhu cầu nước ngoài và tham nhũng.
\item Kết luận: Nội dung nhấn mạnh việc các nước nghèo đang chịu rủi ro mất rừng cao hơn để phục vụ nước giàu, khớp với tiêu đề (xi) Poorer nations at higher risk.
\item Đáp án: xi
\end{itemize}
\item Câu 20: Paragraph H
\begin{itemize}
\item Nội dung chính: Đoạn văn nói về Trung Quốc đang dẫn đầu trong việc "replenishing their forests" (làm đầy lại rừng của họ) thông qua chương trình trồng rừng "Three-North Shelter Forest Program".
\item Kết luận: Hành động trồng lại cây và khôi phục rừng khớp với tiêu đề (viii) Replanting forests.
\item Đáp án: viii
\end{itemize}
\end{itemize}

\cdsection{Phần II: Câu hỏi 21-26 (Summary Completion)}

\cdnoteinline{Yêu cầu: Hoàn thành bản tóm tắt về tác động của phá rừng với giới hạn KHÔNG QUÁ 2 TỪ VÀ/HOẶC MỘT SỐ.}

\cdstep{1. Phân tích Từ khóa và Dự đoán}

Chúng ta cần xác định loại từ cần điền dựa trên ngữ pháp và tìm từ khóa trong bài đọc để định vị thông tin.
\begin{itemize}
\item \textbf{(21)} In Britain, for example, \_\_\_\_\_\_\_ people live in areas at risk of flooding.
\begin{itemize}
\item Từ khóa: "Britain", "people live", "risk of flooding".
\item Dự đoán: Cần một con số chỉ số lượng người.
\item Định vị: Đoạn B, câu đầu tiên: "There are 5 million people living in areas deemed at risk of flooding in England and Wales." (Anh và xứ Wales là đại diện cho Britain).
\end{itemize}
\item \textbf{(22)} Trees catch and \_\_\_\_\_\_\_ water lowering the chance of flooding.
\begin{itemize}
\item Từ khóa: "Trees catch", "water", "lowering chance of flooding".
\item Dự đoán: Cần một động từ mô tả hành động của cây với nước.
\item Định vị: Đoạn B, câu: "The trees catch and store water..."
\end{itemize}
\item \textbf{(23)} By removing trees land erosion is also higher, increasing the chance of \_\_\_\_\_\_\_.
\begin{itemize}
\item Từ khóa: "removing trees", "erosion", "increasing chance of".
\item Dự đoán: Cần một danh từ chỉ hậu quả của xói mòn (erosion).
\item Định vị: Đoạn B, câu: "...increasing the risk of landslides..."
\end{itemize}
\item \textbf{(24)} Deforestation also affects global warming by contributing 15\% of the \_\_\_\_\_\_\_ of greenhouse gasses.
\begin{itemize}
\item Từ khóa: "contributing 15\%", "greenhouse gasses".
\item Dự đoán: Cần một danh từ đi sau "15\% of the...".
\item Định vị: Đoạn C, câu: "Deforestation releases 15\% of all greenhouse gas emissions."
\end{itemize}
\item \textbf{(25)} To make sure that the cutting down of trees is done in a sustainable way, good forestry \_\_\_\_\_\_\_ is important.
\begin{itemize}
\item Từ khóa: "sustainable way", "good forestry", "important".
\item Dự đoán: Cần một danh từ ghép với "forestry".
\item Định vị: Đoạn E, câu đầu tiên: "Forestry management is important to make sure that stocks are not depleted..."
\end{itemize}
\item \textbf{(26)} One country that is reversing this trend is China, making it one of the few nations to \_\_\_\_\_\_\_ the more common negative trend.
\begin{itemize}
\item Từ khóa: "China", "reversing this trend", "few nations to".
\item Dự đoán: Cần một động từ chỉ hành động ngược lại với xu hướng tiêu cực.
\item Định vị: Đoạn H, câu cuối cùng: "...making China one of the few developing nations to reverse the negative trend."
\end{itemize}
\end{itemize}

\cdstep{2. Đối chiếu và Kết luận Đáp án}

Câu

Phân tích và Dẫn chứng

Đáp án

21

Tìm thấy số lượng người sống trong vùng nguy cơ lũ lụt ở Anh tại đoạn B là 5 triệu.

5 million

22

Đoạn B nói cây "catch and store water" (giữ và lưu trữ nước).

store

23

Đoạn B chỉ ra việc xói mòn đất làm tăng nguy cơ sạt lở đất (landslides).

landslides

24

Đoạn C nêu rõ phá rừng thải ra 15\% lượng khí thải (emissions) nhà kính.

emissions

25

Đoạn E nhấn mạnh tầm quan trọng của quản lý (management) lâm nghiệp.

management

26

Đoạn H nói Trung Quốc là một trong số ít quốc gia đảo ngược (reverse) xu hướng tiêu cực.

reverse

\end{document}
